\documentclass[11pt,letterpaper]{article}
\usepackage[margin=1in]{geometry}
\usepackage{amsmath, amssymb}
\usepackage{booktabs}
\usepackage{graphicx}
\usepackage{hyperref}
\usepackage{xcolor}
\usepackage{listings}
\usepackage{caption}
\usepackage{subcaption}

\lstset{
  basicstyle=\ttfamily\small,
  breaklines=true,
  frame=single,
  backgroundcolor=\color{gray!10},
}

\title{Chemical and Electrochemical Stability Analysis\\
       of Fe-Halide Cathodes with Halide Solid Electrolytes\\[0.5em]
       \large Computational Methodology and Results Summary}
\author{Howard Tu}
\date{\today}

\begin{document}
\maketitle
\tableofcontents
\newpage

% ============================================================
\section{Project Overview}
% ============================================================

This report documents the computational analysis of chemical and electrochemical
stability at the interface between iron-halide cathode materials and halide solid
electrolytes (SEs). The study targets three cathode compositions:

\begin{itemize}
  \item \textbf{Row 3}: FeF$_3$ cathode vs.\ five halide SEs
  \item \textbf{Row 4}: FeCl$_3$ cathode vs.\ five halide SEs
  \item \textbf{Row 2}: FeF$_3$+FeCl$_3$ (1:1 molar) cathode vs.\ five halide SEs
\end{itemize}

The five halide solid electrolytes studied are:

\begin{center}
\begin{tabular}{ll}
\toprule
Label & Formula \\
\midrule
Li$_3$YCl$_6$                                   & Li$_3$YCl$_6$ \\
Li$_{2.4}$Y$_{0.4}$Zr$_{0.6}$Cl$_6$               & Li$_{2.4}$Y$_{0.4}$Zr$_{0.6}$Cl$_6$ \\
Li$_3$ScCl$_6$                                  & Li$_3$ScCl$_6$ \\
Li$_{2.375}$Sc$_{0.375}$Zr$_{0.625}$Cl$_6$        & Li$_{2.375}$Sc$_{0.375}$Zr$_{0.625}$Cl$_6$ \\
Li$_3$InCl$_6$                                  & Li$_3$InCl$_6$ \\
\bottomrule
\end{tabular}
\end{center}

All thermodynamic data are retrieved from the Materials Project (MP) database
via the \texttt{mp-api} client and processed using \texttt{pymatgen}.

% ============================================================
\section{Theoretical Framework}
% ============================================================

\subsection{Chemical Stability (Closed System)}

The chemical reaction energy at a mixing ratio $x$ (molar fraction of cathode)
in the closed system is:
\begin{equation}
  \Delta E_{\rm rxn}(x) = E_{\rm hull}^{\rm PD}\!\left[c(x)\right]
    - \left[ x \cdot E_{\rm hull}^{\rm PD}(C) + (1-x) \cdot E_{\rm hull}^{\rm PD}(SE) \right]
  \label{eq:chem}
\end{equation}
where $c(x) = x\,C + (1-x)\,SE$ is the mixed composition, and
$E_{\rm hull}^{\rm PD}$ denotes the energy on the convex hull of the standard
\texttt{PhaseDiagram} (per atom, relative to the stable phases in the
relevant chemical subsystem). A negative $\Delta E_{\rm rxn}$ means the
interface is thermodynamically unstable and will spontaneously react.

\subsection{Electrochemical Stability (Open Li System)}

In a battery, lithium can flow into or out of the interface region.
This is modeled using the \emph{grand potential} phase diagram (GPCD),
in which Li is treated as an open element with chemical potential
\begin{equation}
  \mu_{\rm Li}(V) = \mu_{\rm Li}^{\rm ref} - V
  \label{eq:muli}
\end{equation}
where $V$ is the electrode voltage (V vs.\ Li/Li$^+$) and
$\mu_{\rm Li}^{\rm ref}$ is the Li chemical potential at the most
Li-rich transition point on the phase diagram:
\begin{lstlisting}[language=Python]
mu_li_ref = pd_obj.get_transition_chempots(Element("Li"))[0]
chempots = {"Li": -voltage + mu_li_ref}
gpd = GrandPotentialPhaseDiagram(entries, chempots)
\end{lstlisting}
The grand-potential energy of a composition $c$ is:
\begin{equation}
  \Phi(c) = E_{\rm DFT}(c) - n_{\rm Li}(c)\,\mu_{\rm Li}
\end{equation}
The electrochemical reaction energy is then computed analogously to
Eq.~\eqref{eq:chem} but using $\Phi$ evaluated on the GPCD hull instead
of $E$ on the PD hull.

\paragraph{Physical interpretation.}
At \emph{high voltage} (oxidizing conditions), the Li chemical potential is
very negative ($\mu_{\rm Li} \ll 0$), so it costs energy to reduce Li from
the interface; the cathode and SE are more stable (less negative
$\Delta E_{\rm rxn}$). At \emph{low voltage} (reducing conditions),
Li is readily inserted, driving strong reduction reactions and more negative
$\Delta E_{\rm rxn}$. This is why the low-voltage (blue) curves in the
plots sit at the most negative energies.

\subsection{Voltage Range and Normalization}

\begin{itemize}
  \item Voltage sweep: 1.0--4.2\,V vs.\ Li/Li$^+$, step 0.01\,V (321 steps)
  \item Energy normalization: eV per atom of the reactant \emph{mixture}
        (not per formula unit); this is ``\texttt{norm=True}'' in
        \texttt{InterfacialReactivity}.
  \item $x$ axis: molar fraction of cathode, $x = \frac{n_C/N_C}{n_C/N_C + n_{SE}/N_{SE}}$,
        where $n_C$, $n_{SE}$ are the molar amounts and $N_C$, $N_{SE}$
        are the number of atoms per formula unit.
\end{itemize}

% ============================================================
\section{Rows 3 and 4: Binary Interface Calculations}
% ============================================================

\subsection{Method: \texttt{InterfacialReactivity}}

For the single-cathode cases (FeF$_3$ or FeCl$_3$ vs.\ each SE),
pymatgen's built-in \texttt{InterfacialReactivity} and
\texttt{GrandPotentialInterfacialReactivity} classes are used directly.
These classes scan the tie-line between two compositions in the (grand
potential) phase diagram, locate the ``kink'' points where the set of
stable decomposition products changes, and return the reaction energy
at each kink.

\begin{lstlisting}[language=Python]
# Chemical stability (closed system):
iface = InterfacialReactivity(
    comp1, comp2, pd_obj,
    norm=True, use_hull_energy=False)

# Electrochemical stability (open Li, voltage V):
chempots = {"Li": -voltage + mu_li_ref}
gpd = GrandPotentialPhaseDiagram(entries, chempots)
iface = GrandPotentialInterfacialReactivity(
    comp1, comp2, gpd,
    pd_non_grand=pd_obj,
    include_no_mixing_energy=True,
    norm=True, use_hull_energy=False)

for _, x_atom, energy, rxn, _kJ in iface.get_kinks():
    x_molar = atom_to_molar(x_atom, n1, n2)
    ...
\end{lstlisting}

The atom-fraction output from \texttt{get\_kinks()} is converted to
molar fraction via:
\begin{equation}
  x_{\rm mol} = \frac{x_{\rm atom}/N_C}{x_{\rm atom}/N_C + (1-x_{\rm atom})/N_{SE}}
\end{equation}

\subsection{Results: FeF$_3$ vs.\ Halide SEs (Row 3)}

Outputs: \texttt{chemical\_stability\_FeF3.png},
\texttt{electrochemical\_stability\_FeF3.png},
\texttt{stability\_table\_FeF3.csv} (652 non-zero-energy rows).

Key observations:
\begin{itemize}
  \item FeF$_3$ shows a clear phase boundary at approximately
        $x \approx 0.57$ for the Cl-based SEs, corresponding to a
        mixed Fe--F--Cl--Li product (LiYF$_4$ + LiCl + Fe metal).
  \item Reaction energies at $V=1$\,V range from $-0.65$ to $-1.36$\,eV/atom
        depending on the SE, indicating significant thermodynamic
        instability at low voltages.
  \item At $V=4.2$\,V the energies approach $\sim -0.1$ to $-0.3$\,eV/atom,
        showing improved but non-negligible instability even at high voltage.
  \item The Zr-substituted electrolytes (Li$_{2.4}$Y$_{0.4}$Zr$_{0.6}$Cl$_6$,
        Li$_{2.375}$Sc$_{0.375}$Zr$_{0.625}$Cl$_6$) show additional intermediate
        kinks due to the larger element space (F, Cl, Zr present).
\end{itemize}

\subsection{Results: FeCl$_3$ vs.\ Halide SEs (Row 4)}

Outputs: \texttt{chemical\_stability\_FeCl3.png},
\texttt{electrochemical\_stability\_FeCl3.png},
\texttt{stability\_table\_FeCl3.csv} (282 non-zero-energy rows).

Key observations:
\begin{itemize}
  \item FeCl$_3$ shows a simpler reaction landscape than FeF$_3$:
        in most cases only a single interior kink exists (at $x=1$,
        i.e., pure cathode), corresponding to FeCl$_3$ + 3\,Li $\to$ Fe + 3\,LiCl.
  \item At $V=1$\,V, $\Delta E_{\rm rxn} \approx -1.29$\,eV/atom,
        substantially negative.
  \item The energy decreases nearly linearly with increasing voltage,
        reaching $\approx -0.07$\,eV/atom near 4.2\,V,
        indicating FeCl$_3$ is relatively stable at high operating voltages.
  \item Because F is absent in the FeCl$_3$ system, no LiF or mixed
        fluoride products form, leading to fewer reaction channels.
\end{itemize}

\subsection{Database Note: Li$_3$YCl$_6$}

Li$_3$YCl$_6$ is not stable on the convex hull in the current MP
database (version 2025.09.25). Pymatgen issues a warning and falls
back to the hull energy, which corresponds to the decomposition
$3\,\text{LiCl} + \text{YCl}_3$. This is physically reasonable and
was accepted as-is. Results for the Y-containing SEs therefore reflect
this projected hull energy rather than the energy of a single
Li$_3$YCl$_6$ phase.

% ============================================================
\section{Row 2: Mixed Cathode FeF$_3$+FeCl$_3$ (1:1 Molar)}
% ============================================================

Two approaches were implemented for the mixed cathode.

\subsection{Option A: Pseudo-Compound Approximation}

\subsubsection{Method}

The 1:1 molar mixture of FeF$_3$ and FeCl$_3$ is combined into a
single pseudo-compound:
\begin{equation}
  1\,\text{mol FeF}_3 + 1\,\text{mol FeCl}_3 \;\longrightarrow\; \text{Fe}_2\text{F}_3\text{Cl}_3
\end{equation}
This composition Fe$_2$F$_3$Cl$_3$ is then used as \texttt{comp1} in the
same \texttt{GrandPotentialInterfacialReactivity} pipeline as Rows 3 and 4.
Fe$_2$F$_3$Cl$_3$ is not a stable MP entry; pymatgen projects it onto the
convex hull of the relevant chemical system (Fe, F, Cl, + SE elements + Li).
In the reaction strings it appears as Fe$_2$(ClF)$_3$.

\subsubsection{Results}

Outputs: \texttt{chemical\_stability\_FeF3-FeCl3.png},
\texttt{electrochemical\_stability\_FeF3-FeCl3.png},
\texttt{stability\_table\_FeF3-FeCl3.csv} (687 rows).

\begin{itemize}
  \item A single dominant kink at $x \approx 0.57$ across all SEs,
        similar to the FeF$_3$ case (same element space, same phase
        space projection).
  \item At $V=1$\,V vs.\ Li$_3$YCl$_6$:
        $\Delta E_{\rm rxn} \approx -0.84$\,eV/atom at $x=0.57$
        and $-1.33$\,eV/atom at $x=1$ (pure cathode boundary),
        intermediate between the pure FeF$_3$ ($-0.65$, $-1.36$)
        and FeCl$_3$ ($-1.29$) values.
  \item The mixing of F and Cl in the pseudo-compound leads to
        LiYF$_4$ + LiCl + Fe as the primary products at the
        intermediate kink.
\end{itemize}

\subsubsection{Limitations}

Because Fe$_2$F$_3$Cl$_3$ is treated as a \emph{single entity}, the
reference energy on the cathode side is the hull energy of Fe$_2$F$_3$Cl$_3$
(projected). If FeF$_3$ and FeCl$_3$ can independently and separately
react with the SE (different products at different spatial locations),
Option A does not capture this. Nor does it capture any intrinsic
chemical reactivity \emph{between} FeF$_3$ and FeCl$_3$ themselves.

% -------------------------------------------------------
\subsection{Option C: Ternary Grand-Potential Approach}
% -------------------------------------------------------

\subsubsection{Motivation}

In reality the cathode is a \emph{physical mixture} of two separate
phases, not a single compound. Option C treats FeF$_3$ and FeCl$_3$
as independent reactants, so all three species---FeF$_3$, FeCl$_3$,
and SE---can participate simultaneously in the reaction at any mixing ratio.

\subsubsection{Composition Parameterization}

The interface is parameterized by the cathode molar fraction $x$:
\begin{equation}
  c(x) = x\!\cdot\!\tfrac{1}{2}\,\text{FeF}_3
        + x\!\cdot\!\tfrac{1}{2}\,\text{FeCl}_3
        + (1-x)\!\cdot\!\text{SE}
  \label{eq:ternary_comp}
\end{equation}
which traces a line in the ternary (or higher-dimensional) composition
space from pure SE at $x=0$ to the 1:1 cathode mixture at $x=1$.

\subsubsection{Reaction Energy Formula}

The grand-potential reaction energy per total atom of the mixture is:
\begin{equation}
  \Delta E_{\rm rxn}(x) =
    \Phi_{\rm hull}\!\left[c(x)\right]
    - \left[
        \alpha_A\,\Phi_{\rm hull}(\text{FeF}_3)
      + \alpha_B\,\Phi_{\rm hull}(\text{FeCl}_3)
      + \alpha_{SE}\,\Phi_{\rm hull}(\text{SE})
      \right]
  \label{eq:ternary_erxn}
\end{equation}
where $\alpha_i = N_{\rm atoms}^{(i)} / N_{\rm atoms}^{\rm total}$ is
the atom fraction of component $i$ in the mixture (counting all atoms
including Li).

\subsubsection{Grand-Potential Phase Diagram and the Li Projection}

The GPCD is built in the \emph{Li-free} composition subspace: Li is the
open element, so pymatgen removes it from all entries and builds the hull
over the remaining elements (e.g., \{Fe, F, Cl, Y\} for the
FeF$_3$/Li$_3$YCl$_6$ system). This means:

\begin{enumerate}
  \item The GPCD's \texttt{get\_hull\_energy\_per\_atom(comp)} method
        \emph{does not accept Li-containing compositions directly}.
        Passing Li$_3$YCl$_6$ raises a \texttt{ValueError} because Li
        is not in the diagram's element list.
  \item All compositions must have Li stripped before querying the GPCD.
  \item The hull energy returned is \emph{per non-Li atom}.
\end{enumerate}

\paragraph{Normalization fix.}
To keep all energies on a consistent ``per total atom (including Li)''
basis, the following rescaling is applied:
\begin{equation}
  \Phi_{\rm hull}^{\rm per\,total}(c)
  = \Phi_{\rm hull}^{\rm per\,non\text{-}Li}(\tilde{c})
    \cdot \frac{N_{\rm non\text{-}Li}}{N_{\rm total}}
  \label{eq:rescale}
\end{equation}
where $\tilde{c}$ is the composition with Li removed and $N_{\rm non\text{-}Li}$,
$N_{\rm total}$ are the corresponding atom counts.

\paragraph{Why this is self-consistent.}
Consider Li$_3$YCl$_6$ ($N_{\rm total}=10$, $N_{\rm non\text{-}Li}=7$,
$\tilde{c}=\text{YCl}_6$):
\begin{equation*}
  \Phi_{\rm hull}^{\rm per\,total}(\text{Li}_3\text{YCl}_6)
  = \Phi_{\rm hull}^{\rm per\,non\text{-}Li}(\text{YCl}_6) \cdot \tfrac{7}{10}
  = \frac{E(\text{Li}_3\text{YCl}_6) - 3\mu_{\rm Li}}{10}
\end{equation*}
which is exactly the grand-potential energy per total formula-unit atom.
At $x=0$ (pure SE): $\Phi_{\rm hull}^{\rm per\,total}(c_{\rm SE}) = \Phi_{\rm hull}^{\rm per\,total}(\text{SE})$,
so the reference and product energies cancel and $\Delta E_{\rm rxn}(0)=0$ as required.
For FeF$_3$ and FeCl$_3$ (no Li): $N_{\rm non\text{-}Li}=N_{\rm total}$,
so Eq.~\eqref{eq:rescale} reduces to the standard GPCD query.

\subsubsection{Dense Scan and Kink Detection}

Because there is no built-in ternary \texttt{InterfacialReactivity} class,
the reaction energy is computed by a dense scan:
\begin{lstlisting}[language=Python]
xs = np.linspace(0, 1, N_SCAN=300)  # molar fraction grid
for x in xs:
    comp_mix = x/2*FeF3 + x/2*FeCl3 + (1-x)*SE  # element-wise
    e_hull   = gpd_hull_per_total_atom(comp_mix, gpd)
    e_ref    = (a_A*e_A + a_B*e_B + a_SE*e_SE) / total_atoms
    E_rxn[x] = e_hull - e_ref
\end{lstlisting}
Kink positions (phase boundaries) are then detected by finding points
where the slope $|dE_{\rm rxn}/dx|$ changes abruptly
(threshold \texttt{SLOPE\_TOL}=$10^{-3}$\,eV/atom per molar fraction unit).

\subsubsection{Results}

Outputs: \texttt{chemical\_ternary\_all.png},
\texttt{electrochemical\_ternary\_all.png},
\texttt{ternary\_table.csv} (1346 rows).

Key differences from Option A:
\begin{itemize}
  \item \textbf{Multiple kinks} appear between $x \approx 0.73$ and
        $x = 1.0$ that are absent in Option A. These reflect a richer
        ternary phase boundary structure when FeF$_3$ and FeCl$_3$ are
        allowed to react independently and with each other.
  \item \textbf{Smaller magnitude} near the cathode-rich end
        ($\Delta E_{\rm rxn} \approx -0.03$ to $-0.04$\,eV/atom at
        $x \approx 0.73$, $V=1$\,V for Li$_3$YCl$_6$) compared to
        Option A's single kink value of $-0.84$\,eV/atom at $x=0.57$.
        This is because at large $x$ (little SE present) the reaction
        is mainly between the small amount of SE and the large excess
        of cathode---a dilute perturbation.
  \item The dense scan produces continuous curves with visible slope
        changes (kinks), whereas Option A produces a piecewise-linear
        curve with exactly the kink points returned by
        \texttt{get\_kinks()}.
  \item At $x=1$ (pure cathode, no SE), $\Delta E_{\rm rxn}$ may be
        slightly non-zero, reflecting any thermodynamic driving force
        for FeF$_3$ and FeCl$_3$ to react with each other (form mixed
        Fe--F--Cl phases). This quantity is absent in Option A by
        construction.
\end{itemize}

\subsubsection{Comparison of Option A vs.\ Option C}

\begin{center}
\begin{tabular}{lll}
\toprule
                         & \textbf{Option A}             & \textbf{Option C} \\
\midrule
Cathode representation   & Single pseudo-compound        & Two separate phases \\
                         & Fe$_2$F$_3$Cl$_3$                 & FeF$_3$ + FeCl$_3$ \\
Reference energy (cathode) & Hull energy of Fe$_2$F$_3$Cl$_3$ & $\frac{1}{2}E_{\rm hull}(\text{FeF}_3)
                                                            +\frac{1}{2}E_{\rm hull}(\text{FeCl}_3)$ \\
Composition space        & Binary (cathode vs.\ SE)      & Ternary (FeF$_3$, FeCl$_3$, SE) \\
Kinks found              & 1--2 per voltage              & Multiple per voltage \\
FeF$_3\leftrightarrow$FeCl$_3$ reaction & Absorbed & Explicitly captured \\
Method used              & \texttt{GrandPotentialInterfacialReactivity} & Dense scan + GPCD hull \\
\bottomrule
\end{tabular}
\end{center}

% ============================================================
\section{Output Files Summary}
% ============================================================

\begin{center}
\begin{tabular}{lll}
\toprule
File & Row & Description \\
\midrule
\texttt{chemical\_stability\_FeF3.png}         & 3 & Chemical, all 5 SEs \\
\texttt{electrochemical\_stability\_FeF3.png}  & 3 & Electrochemical, 5 subplots \\
\texttt{stability\_table\_FeF3.csv}            & 3 & 652 rows, 0.1\,V steps \\
\texttt{chemical\_stability\_FeCl3.png}        & 4 & Chemical, all 5 SEs \\
\texttt{electrochemical\_stability\_FeCl3.png} & 4 & Electrochemical, 5 subplots \\
\texttt{stability\_table\_FeCl3.csv}           & 4 & 282 rows, 0.1\,V steps \\
\texttt{chemical\_stability\_FeF3-FeCl3.png}        & 2A & Chemical, pseudo-compound \\
\texttt{electrochemical\_stability\_FeF3-FeCl3.png} & 2A & Electrochemical, pseudo-compound \\
\texttt{stability\_table\_FeF3-FeCl3.csv}           & 2A & 687 rows, 0.1\,V steps \\
\texttt{chemical\_ternary\_all.png}            & 2C & Chemical, ternary method \\
\texttt{electrochemical\_ternary\_all.png}     & 2C & Electrochemical, ternary method \\
\texttt{ternary\_table.csv}                    & 2C & 1346 rows, 0.1\,V steps \\
\bottomrule
\end{tabular}
\end{center}

% ============================================================
\section{Software and Data Sources}
% ============================================================

\begin{itemize}
  \item \textbf{Database}: Materials Project (MP), accessed via \texttt{mp-api}
        client; database version 2025.09.25.
  \item \textbf{pymatgen} version: 2025.10.7 (updated API; \texttt{pd\_non\_grand}
        is now a required positional argument in
        \texttt{GrandPotentialInterfacialReactivity}).
  \item \textbf{DFT functional}: GGA (PBE) for all Fe-halide entries.
        No GGA+U correction is applied to Fe in halide environments
        within the current MP database.
  \item \textbf{Main scripts}:
    \begin{itemize}
      \item \texttt{electrochemical\_stability.py} --- Rows 2A, 3, 4 via
            \texttt{InterfacialReactivity}
      \item \texttt{ternary\_stability.py} --- Row 2C via dense GPCD scan
    \end{itemize}
\end{itemize}

\end{document}
